\subsubsection{Matrices de loci chromosomiques à nombre de cœurs fixé}\label{sec:matrice-chromosome}
La comparaison sur une même matrice sépare l'effet du noyau de celui du nombre de cœurs. Le chromosome~1 de CHM13~v1.1 fournit 768 fenêtres disjointes. Chacune des quatre longueurs, entre 512 et 3\,072 bases, compte 192 fenêtres. Ces loci éloignés constituent une charge de recherche entre régions génomiques. La cohorte d'haplotypes de la section~\ref{sec:cohortecompacte} porte sur une même région chez plusieurs personnes.

Un rang peut construire une seule fois l'index de dernière occurrence de chaque cible et le partager entre huit travailleurs. Chaque travailleur conserve deux lignes de scores pour ses paires successives. L'index permet aussi de reconnaître une distance minimale avant le calcul. Pour \(m\le n\), toute transformation coûte au moins \(n-m\) éditions. Elle atteint cette borne si et seulement si \(A\) est une sous-séquence de \(B\). Les \(n-m\) insertions nécessaires épuisent alors le budget. L'algorithme~\ref{alg:subseq-chr} essaie la borne lorsque \(n\ge 3m\). Ce seuil règle le coût de la recherche.

\begin{algorithm}[!htbp]
\caption{Reconnaissance d’une sous-séquence par dernières occurrences}\label{alg:subseq-chr}
\begin{algorithmic}[1]
\Require Strings \(A[1:m], B[1:n]\), \(m\le n\), last-occurrence index \(K_B\), score solver \(F\)
\Ensure Global edit distance \(d(A,B)\)
\If{\(n\ge 3m\)}
  \State \(p\gets n\)
  \For{\(i=m,m-1,\ldots,1\)}
    \State \(k\gets K_B[A_i,p]\)
    \If{\(k=0\)} \State \Return \(F(A,B)\) \EndIf
    \State \(p\gets k-1\)
  \EndFor
  \State \Return \(n-m\)
\EndIf
\State \Return \(F(A,B)\)
\end{algorithmic}
\end{algorithm}

La figure~\ref{fig:matrice-chromosome} compare les solveurs dans cinq blocs permutés. Le chronomètre inclut la construction des index et la réduction MPI. Avec \(W_q=\sum_{i<j}\ell_i\ell_j\) pour \(q\) loci, les efficacités entre un et deux nœuds s'écrivent
\[
  E_{\mathrm s}=\frac{T_{768,1}}{2T_{768,2}},\qquad
  E_{\mathrm w}=\frac{W_{768}}{2W_{544}}\frac{T_{544,1}}{T_{768,2}}.
\]
Le rapport \(W_{768}/(2W_{544})\) vaut 0,9973. Pour Edlib et WFA2, \(W_q\) décrit la surface du problème indépendamment du nombre de cellules visitées.
\chart[.95\linewidth]{distributed/matrice_chromosome.pdf}{Temps et efficacités sur EPYC pour les mêmes distances globales. Les intervalles fins donnent l'IC de Student à 95\,\% et les traits épais un écart-type sur cinq blocs. Les barres hachurées incluent le test de sous-séquence.}{fig:matrice-chromosome}{bench/mpi/chromosome/measure.py}

Sur 384 cœurs, le partage des index et des tampons ramène NvPD de \(1{,}7449\pm0{,}0053\) à \(1{,}3402\pm0{,}0035\)\,s en ordre de Hilbert pondéré. Avec les mêmes index, le vol de travail prend \(1{,}1679\pm0{,}0039\)\,s, contre \(1{,}4777\pm0{,}0041\)\,s par blocs statiques. Le produit des longueurs prédit le nombre de cellules denses. L'ordre d'accès aux cibles et le coût de la file interviennent aussi dans le temps mesuré. L'efficacité forte du calcul indexé est de 75,0\,\% et son efficacité à charge par cœur fixée de 75,6\,\%.

Le test de sous-séquence résout 39\,679 paires parmi les 110\,592 examinées. Avec WFA2, il ramène la durée de \(1{,}2931\pm0{,}0021\) à \(1{,}0742\pm0{,}0053\)\,s. Le solveur bit-parallèle d'Edlib traite déjà plusieurs états d'une colonne dans un mot machine. Son temps passe de 0,1680 à 0,1645\,s.

\synthese{Sur les 294\,528 paires, Edlib termine en \(0{,}1680\pm0{,}0006\)\,s sur 384 cœurs. L'écart restant avec NvPD indexé atteint un facteur de 7,98. Les comparaisons suivantes portent sur des régions homologues dont la distance est plus faible.}
